% Laget automatisk av elevark/lag_elevark.py fra sim/elektrisk-felt.html. Endringer her blir overskrevet.
\documentclass[11pt]{article}
\usepackage{elevark}
\begin{document}
\elevarkhode{Fysikk 2 · Elektriske felt og krefter}{Elektrisk felt}{https://fysikklab.no/sim/elektrisk-felt.html}{Simuleringen}

Rundt en ladning er det et elektrisk felt. Dra ladningene og målepunktet $P$ rundt, og se hvordan feltet endrer seg. Når det er flere ladninger, legger vi sammen feltene fra hver ladning som vektorer.

\section*{Oppgaver}

\begin{oppgaveliste}
\oppgave \fase{gjett}{Gjett}Velg «Én ladning». Flytt $P$ slik at den er 5,0~cm fra ladningen, og så 10,0~cm fra ladningen. Hvor mange ganger svakere tror du feltet blir? Test.
\svar
\oppgave \fase{test}{Regn}Ladningen er +2,0~nC. Regn ut feltstyrken 10,0~cm fra ladningen, og sjekk med simuleringen. Hvor stor kraft virker på et elektron i dette punktet, og i hvilken retning?
\svar
\oppgave \fase{forklar}{Forklar}Velg «Pluss og minus» og slå på feltpilene. Hvor er feltet sterkest? Hvordan ser du det på feltlinjene? Hvorfor kan to feltlinjer aldri krysse hverandre?
\svar
\oppgave \fase{test}{Regn}Ladningene i «Pluss og minus» er +2,0~nC og −2,0~nC, med 12~cm mellom. Regn ut kraften mellom dem med Coulombs lov. Regn så ut feltstyrken midt mellom dem, og sjekk svarene. Slå på «Krefter på ladningene». Hva viser pilene?
\svar
\oppgave \fase{gjett}{Gjett}Velg «+4~nC og +1~nC». Et sted mellom ladningene er feltet null. Ligger punktet nærmest den store eller den lille ladningen? Regn ut hvor det er, og finn det med $P$.
\svar
\oppgave \fase{test}{Regn}Velg «To like ladninger». Hva er feltet midt mellom ladningene? $P$ ligger 5,5~cm over midtpunktet. Regn ut feltet i $P$ ved å dekomponere feltet fra hver ladning. Sjekk svaret.
\svar
\oppgave \fase{utfordring}{Utfordring}Velg «Tre ladninger». Ladningene ligger i hjørnene av en likesidet trekant med sider på 12~cm. Hva er retningen på feltet i $P$? Lag selv et oppsett med fire ladninger der feltet er null midt i, og forklar hvorfor.
\svar
\end{oppgaveliste}

\end{document}
